% TODO make hw stylesheet
\documentclass[10pt]{article}
\usepackage[letterpaper,top=1in]{geometry}
\usepackage{quiver}
\usepackage[dvipsnames,svgnames]{xcolor}
\usepackage[shortlabels]{enumitem}
\usepackage[framemethod=TikZ]{mdframed}
\usepackage{amsmath,amssymb,amsthm}
\usepackage[colorlinks]{hyperref}
\usepackage{multicol}
\usepackage{tabularx}
\usepackage{stmaryrd}

\hypersetup{
	colorlinks=true,
	linkcolor=blue,
	filecolor=magenta,
	urlcolor=magenta,
	pdftitle={MAT 534 HW}
}

\usepackage{mathtools}
\usepackage{thmtools}
\usepackage[textsize=footnotesize]{todonotes}
\usepackage{titling}

\reversemarginpar
\mdfdefinestyle{mdbluebox}{roundcorner=10pt,innerbottommargin=9pt,
	linecolor=blue,backgroundcolor=TealBlue!5,}
\declaretheoremstyle[headfont=\sffamily\bfseries\color{MidnightBlue},
mdframed={style=mdbluebox},]{thmbluebox}

\mdfdefinestyle{mdbloobox}{frametitlefont=\bfseries,innerbottommargin=8pt,
	nobreak=true,backgroundcolor=TealBlue!5,linecolor=blue,}
\declaretheoremstyle[headfont=\bfseries\color{MidnightBlue},
mdframed={style=mdbloobox},headpunct={\\[3pt]},postheadspace=0pt,]{thmbloobox}

\mdfdefinestyle{mdredbox}{frametitlefont=\bfseries,innerbottommargin=8pt,
	nobreak=true,backgroundcolor=Salmon!5,linecolor=RawSienna,}
\declaretheoremstyle[headfont=\bfseries\color{RawSienna},
mdframed={style=mdredbox},headpunct={\\[3pt]},postheadspace=0pt,]{thmredbox}

\mdfdefinestyle{mdgreenbox}{linecolor=ForestGreen,backgroundcolor=ForestGreen!5,
	linewidth=2pt,rightline=false,leftline=true,topline=false,bottomline=false,}
\declaretheoremstyle[headfont=\bfseries\sffamily\color{ForestGreen!70!black},
mdframed={style=mdgreenbox},headpunct={ --- },]{thmgreenbox}

\mdfdefinestyle{mdorangebox}{linecolor=RedOrange,backgroundcolor=RedOrange!5,
	linewidth=2pt,rightline=false,leftline=true,topline=false,bottomline=false,}
\declaretheoremstyle[headfont=\bfseries\sffamily\color{RedOrange!70!black},
mdframed={style=mdorangebox},headpunct={ --- },]{thmorangebox}

\mdfdefinestyle{mdblackbox}{linecolor=black,backgroundcolor=RedViolet!5!gray!5,
	linewidth=3pt,rightline=false,leftline=true,topline=false,bottomline=false,}
\declaretheoremstyle[mdframed={style=mdblackbox}]{thmblackbox}

\declaretheoremstyle[spaceabove=6pt,spacebelow=6pt]{basehead}

\declaretheorem[style=basehead,name=Problem]{problem}
\declaretheorem[style=thmredbox,name=Problem,numbered=no]{problem*}
\declaretheorem[style=thmbloobox,name=Setup]{setup} % for config geo
\declaretheorem[style=thmredbox,name=Example,sibling=problem]{example}
\declaretheorem[style=thmredbox,name=$\bigstar$ Problem,sibling=problem]{niceprob}
\declaretheorem[style=thmbluebox,name=Theorem,sibling=problem]{theorem}
\declaretheorem[style=thmbluebox,name=Corollary,sibling=theorem]{corollary}
\declaretheorem[style=thmbluebox,name=Proposition,sibling=theorem]{prop}
\declaretheorem[style=thmbluebox,name=Lemma,numberwithin=problem]{lemma}
\declaretheorem[style=thmbluebox,name=Theorem,numbered=no]{theorem*}
\declaretheorem[style=thmbluebox,name=Lemma,numbered=no]{lemma*}
\declaretheorem[style=thmgreenbox,name=Claim,numberwithin=problem]{claim}
\declaretheorem[style=thmgreenbox,name=Claim,numbered=no]{claim*}
\declaretheorem[style=thmblackbox,name=Remark,numberwithin=problem]{remark}
\declaretheorem[style=thmblackbox,name=Remark,numbered=no]{remark*}
\declaretheorem[style=thmgreenbox,name=Definition,sibling=theorem]{definition}
\declaretheorem[style=thmgreenbox,name=Definition,numbered=no]{definition*}

\providecommand{\ol}{\overline}
\providecommand{\quiv}{\equiv}
\providecommand{\ceil}[1]{\left\lceil #1 \right\rceil}
\providecommand{\floor}[1]{\left\lfloor #1 \right\rfloor}
\newcommand{\dd}{\mathrm{d}}
\providecommand{\ul}{\underline}
\providecommand{\eps}{\varepsilon}
\providecommand{\half}{\frac{1}{2}}
\providecommand{\CC}{\mathbb C}
\providecommand{\FF}{\mathbb F}
\providecommand{\NN}{\mathbb N}
\providecommand{\QQ}{\mathbb Q}
\providecommand{\RR}{\mathbb R}
\providecommand{\ZZ}{\mathbb Z}
\providecommand{\dg}{^\circ}
\providecommand{\ii}{\item}
\providecommand{\setdiff}{\smallsetminus}
\providecommand{\alert}[1]{{\sffamily\textbf{\textcolor{blue}{#1}}}}
\providecommand{\inv}{^{-1}}
\providecommand{\mb}[1]{\mathbf{#1}}

\newenvironment{soln}{\begin{proof}[Solution]}{\end{proof}}

\providecommand{\pow}{\operatorname{Pow}}
\providecommand{\aut}{\operatorname{Aut}}
\providecommand{\ann}{\operatorname{Ann}}
\providecommand{\vocab}[1]{{\textbf{\textcolor{ForestGreen}{#1}}}}
\providecommand{\lcm}{\operatorname{lcm}}
\providecommand{\abs}[1]{\left|#1\right|}
\providecommand{\Ob}{\operatorname{Ob}}
\providecommand{\Hom}{\operatorname{Hom}}
\providecommand{\Aut}{\operatorname{Aut}}
\providecommand{\id}{\operatorname{id}}
\newcommand{\doubleparen}[1]{(\!(#1)\!)}

\usepackage{mathrsfs}

% place title at top right
\setlength{\droptitle}{-8ex}
\pretitle{\begin{flushright}\Large\bfseries}
\posttitle{\par\end{flushright}}
\preauthor{\begin{flushright}}
\postauthor{\end{flushright}}
\predate{\begin{flushright}}
\postdate{\end{flushright}}

\title{MAT 534 HW09}
\author{\large Aditya Pahuja}
\date{\large Due 11/12}
\begin{document}
\maketitle

\begin{problem}
    A module $M$ is said to be \vocab{simple} if
    $M$ is not the zero module and the only submodules of $M$
    are $M$ and $0$. Show that $M$ is a simple module
    if and only if $M\cong R/I$, where $I$ is a maximal ideal of $R$
    (where isomorphism is as an $R$-module).
\end{problem}

\begin{soln}
    Let $m$ be any nonzero element of simple module $M$. Then
    $\{rm : r\in R\}$ is a nontrivial submodule of $M$, so it equals $M$.
    Let $I$ be the annihilator of $m$, which by Problem 4(a) is an ideal of $R$;
    then, by (my solution to) Problem 4(b), $M\cong R/I$ with isomorphism
    $\phi_m\colon R/I\to M$ given by $\phi_m(r + I) = rm$.
    This means that given nonzero $s + I\in R/I$,
    there exists $r + I\in R/I$ such that
    $\phi_{sm}(r + I) = rsm = m$, i.e. $rs - 1\in\ann(m) = I$,
    so $r+I$ and $s+I$ are inverse in $R/I$.
    Thus, $R/I$ is a field, making $I$ a maximal ideal.

    Conversely, we show that $R/I$ is a simple $R$-module
    if $I$ is a maximal ideal of $R$.
    Since $R/I$ is a field, multiplication by
    a nonzero element of $R/I$ is a bijection on $R/I$,
    meaning that given $r+I\in R/I$, the set $R(r + I) = \{s(r + I) : s\in R\}$
    is equal to $R/I$. In particular, if $N$ is a nontrivial submodule
    of $R/I$ containing a nonzero $m$, then $R/I = Rm \subseteq N$,
    so $N$ is the entire module.
\end{soln}

\begin{problem}
    Suppose that $M$ and $N$ are submodules of a module $P$.
    Show that $(M+N)/M\cap N\cong (M/M\cap N)\oplus(N/M\cap N)$.
\end{problem}

\begin{soln}
    Let $\phi\colon (M/M\cap N)\oplus(N/M\cap N)\to (M+N)/M\cap N$
    such that
    \begin{equation*}
        \phi(a+M\cap N,b+M\cap N) = (a+b) + M\cap N.
    \end{equation*}
    This is of course surjective because any element of $M+N$
    is written as $a+b$ for some $a\in M$, $b\in N$.
    To show it is injective, we need to show that if
    $a+b\quiv c+d\pmod{M\cap N}$ for $a,c\in M$ and $b,d\in N$,
    then $a\quiv c\pmod{M\cap N}$ and $b\quiv d\pmod{M\cap N}$.
    Indeed, we know a priori that $(a-c) + (b-d) \in M\cap N$,
    so since $a-c\in M$ and $(a-c) + (b-d)\in M$,
    their difference $b-d$ is also in $M$, hence in $M\cap N$,
    and similarly for $a-c$.
    Thus $\phi$ is an isomorphism between the two modules.
\end{soln}

\begin{problem}
    Suppose that $R$ is an integral domain.
    An element $m\in M$ is said to be a \vocab{torsion element}
    if $rm = 0$ for some $r\ne 0$ in $R$, and we say that
    $r$ \vocab{annihilates} $m$.
    \begin{enumerate}[(a)]
	\ii Let $M_{\text{tor}}$ be the set of torsion elements of $M$.
	Show that $M_{\text{tor}}$ is a submodule of $M$.
	\ii $M$ is said to be \vocab{torsion-free} if $M_{\text{tor}} = 0$.
	Show that $M/M_{\text{tor}}$ is always torsion-free.
	\ii Suppose $R = \ZZ$ and $M = \RR/\ZZ$.
	What is $M_{\text{tor}}$ in this case?
    \end{enumerate}
\end{problem}

\begin{soln}
    (a) Let $m_1,m_2\in M_{\text{tor}}$, so there exist
    $r_1,r_2\in R$ such that $r_1m_1 = r_2m_2 = 0$.
    Then for every $r\in R$, $r_1(rm_1) = r(r_1m_1) = 0$,
    so $M_{\text{tor}}$ is closed under multiplication by $R$,
    and $r_1r_2(m_1 + m_2) = 0$, so $M_{\text{tor}}$
    is closed under addition.
    \\

    (b) Let $m + M_{\text{tor}}$ be a torsion element of $M/M_{\text{tor}}$.
    Then $rm \in M_{\text{tor}}$ for some $r\in R$,
    so there exists $r'\in R$ such that $r'(rm) =(r'r)m = 0$,
    which in turn implies that $m\in M_{\text{tor}}$,
    i.e. $m + M_{\text{tor}}$ is the zero element of $M/M_{\text{tor}}$.
    \\

    (c) The torsion submodule of $\RR/\ZZ$ is $\QQ/\ZZ$,
    since the torsion elements are the cosets corresponding to
    the elements of $\RR$ which, when multiplied by some integer,
    are themselves integers.
\end{soln}

\begin{problem}
    Let $M$ be a module and let $m\in M$.
    The \vocab{annihilator} $\operatorname{Ann}(m)$ of $m$ in $R$
    is the set $\{r\in R : rm = 0\}$.
    \begin{enumerate}[(a)]
        \ii Show that $\ann(m)$ is an ideal of $R$.
	\ii A module $M$ is said to be \vocab{cyclic}
	if it is generated by a single element $m_0\in M$.
	Show that $M$ is cyclic if and only if $M\cong R/I$
	for some ideal $I\subseteq R$.
    \end{enumerate}
\end{problem}

\begin{soln}
    (a) If $r_1,r_2\in\ann(m)$ and $r\in R$ then
    \begin{equation*}
	(r_1+r_2)m = r_1m + r_2m = 0,\qquad
	r(r_1m) = (rr_1)m = 0
    \end{equation*}
    implies that $\ann(m)$ is closed under addition
    and under multiplication by elements of $R$, hence is an ideal of $R$.
    \ii

    (b) Suppose $M$ is generated by $m_0\in M$
    and let $I = \ann(m_0)$. Then, define $\phi\colon R/I\to M$
    by $\phi(r + I) = rm_0$. This map is well-defined because
    if $r-s\in I$ then $rm_0 - sm_0 = (r-s)m_0 = 0$.
    It is surjective because $M$ is generated by $m_0$,
    and it is an injective map because
    \begin{equation*}
        \phi(r + I) = rm_0 = 0\implies r\in\ann(m) = I.
    \end{equation*}
    Thus $\phi$ is an isomorphism between $R/I$ and $M$.
\end{soln}

\begin{problem}
    Suppose that $R$ is an integral domain,
    and let $I$ be a nonzero ideal of $R$.
    Show that $I\cong R$ if and only if $I$ is a principal ideal of $R$.
\end{problem}

\begin{soln}
    Suppose first that $I\cong R$ and let $\phi\colon R\to I$ be an isomorphism.
    Define $a = \phi(1)$. Then, any element of $I$ can be written as
    $\phi(r) = r\phi(1) = ra$ for some $r\in R$, so $I = (a)$.

    Conversely, if $I = (a)$ for some $a\in R$, then the morphism
    $\phi\colon R\to I$ determined by setting $\phi(1) = a$
    is an isomorphism: it is clearly surjective because
    $I = \{ra = \phi(r) : r\in R\}$, and it is injective because
    $ra = r'a$ for $r,r'\in R$ implies $r = r'$, since $R$ is an integral domain.
\end{soln}

\begin{problem}
    Let $F$ be a field and let $V$ be an $F[x]$-module
    that is finite-dimensional as a vector space over $F$.
    \begin{enumerate}[(a)]
        \ii Show that $V$ is a torsion module.
	\ii Let $T\colon V\to V$ be the map defined by
	$T(v) = xv$ for $v\in V$. Show that $T$ is a linear transformation
	of the $F$-vector space $V$.
    \end{enumerate}
\end{problem}

\begin{soln}
    (a) Since $V$ is a finite-dimensional $F$-vector space,
    for each $v\in V$ there exists an $n$ such that
    $v$, $xv$, \dots, $x^nv$ are $F$-linearly independent,
    so there exist $a_0,\dots,a_n\in F$, not all zero, such that
    \begin{equation*}
        a_0v + a_1xv + \cdots + a_nx^nv =
	(a_0 + a_1x + \cdots + a_nx^n)v = 0,
    \end{equation*}
    which implies that $v$ is a torsion element of $V$.
    \\

    (b) $T$ essentially inherits linearity from the fact that
    $V$ is an $F[x]$-module; for $v,w\in V$ and $c\in F$,
    $T(v+w) = x(v+w) = xv + xw = T(v)+T(w)$,
    and $T(cv) = x\cdot cv = c\cdot xv = cT(v)$.
\end{soln}

\pagebreak

\begin{problem}
    Let $F$ be a field, $G$ a finite group, and $F[G]$ the group ring.
    Let $V$ be a finitely generated $F[G]$-module.
    \begin{enumerate}[(a)]
        \ii Show that $V$ is a finite-dimensional vector space over $F$.
	\ii Let $g\in G$, and let $T_g\colon V\to V$ be the map $T_g(v) = gv$.
	Show that $T_g$ is a linear transformation of
	the $F$-vector space $V$.
	\ii Let $v_1$, \dots, $v_n$ be a basis for $V$ over $F$,
	and for each $g\in G$, let $M_g$ be the matrix of $T_g$
	with respect to the basis $v_1$, \dots, $v_n$.
	Show that $M_g\in GL_n(F)$ and that the map
	\begin{equation*}
	    \rho\colon G\to GL_n(F),\qquad g\mapsto M_g
	\end{equation*}
	is a group homomorphism.
    \end{enumerate}
\end{problem}

\begin{soln}
    (a) Obviously $V$ is a vector space over $F$;
    the harder part of the problem is to show that this vector space
    is of finite dimension. Suppose $V$ is generated by $v_1,\dots,v_m\in V$.
    Any element of $V$ is of the form $r_1v_1 + \cdots + r_mv_m$
    where the $r_i$ belong to $F[G]$, meaning $r_i = \sum_{g\in G} a_{i,g}g$.
    This means that any element of $V$ is an $F$-linear combination
    of the elements of the finite set $\{gv_i : g\in G, i=1,\dots,m\}$
    (where the coefficients are the $a_{i,g}$),
    so $V$ is spanned by a finite set, hence has finite dimension.
    \\

    (b) As in the previous problem, linearity is pretty much trivial:
    For $v,w\in V$ and $c\in F$,
    $T_g(v+w) = g(v+w) = gv + gw = T_g(v)+T_g(w)$
    and $T_g(cv) = g\cdot cv = c\cdot gv = cT_g(v)$.
    \\

    (c) To show $M_g\in GL_n(F)$ we need to show that $T_g$ is invertible;
    the inverse is simply $T_{g\inv}$, as
    $T_g(T_{g\inv}(v)) = gg\inv v = v$ and $T_{g\inv}(T_g(v)) = g\inv gv = v$.
    Then, to show that $\rho$ is a homomorphism, we need to show that
    $T_{gh} = T_g\circ T_h$ for any $g,h\in G$; this is true because
    \begin{equation*}
	T_{gh}(v) = ghv = T_g(hv) = T_g(T_h(v)) = (T_g\circ T_h)(v).\qedhere
    \end{equation*}
\end{soln}

\begin{problem}
    Let $R$ be a noetherian ring.
    \begin{enumerate}[(a)]
        \ii Let $M$ be a submodule of $R^n$ for some $n\ge 1$.
	Let $\pi_1\colon R^n\to R$ be the projection to the first coordinate,
	so $\pi_1(r_1,\dots,r_n) = r_1$. Show that $\pi_1(M)$
	is an ideal in $R$.
	\ii Let $M_0 = \{m\in M : \pi_1(m) = 0\}$.
	Show that there is a finitely generated submodule $M'\subseteq M$
	such that $M = M' + M_0$.
	\ii Prove that every submodule of $R^n$ is finitely generated.
	\ii Deduce that every submodule of a finitely generated $R$-module
	is again finitely generated.
    \end{enumerate}
\end{problem}

\begin{soln}
    (a) $\pi_1(M)$ is a submodule of the $R$-module $R$,
    and submodules of rings are ideals.
    \\

    (b) Since $R$ is noetherian, $\pi_1(M)$ is finitely generated,
    say by $r_1,\dots,r_k\in\pi_1(M)$.
    Let $m_1,\dots,m_k\in M$ be such that $\pi_1(m_j) = r_j$
    for $1\le j\le k$, and let $M'$ be the submodule generated by the $m_j$.
    Then, $\pi_1(M') = \pi_1(M)$ by construction,
    so if $m\in M$, there exists $m'\in M'$ such that
    the first coordinate of $m$ is equal to that of $m'$.
    This means $\pi_1(m-m') = 0$, so $m-m'\in M_0$,
    proving that $M' + M_0 = M$.
    \\

    (c) We induct on $n$, with the $n=1$ case following from
    the fact that every submodule (ideal) of $R$ is finitely generated
    because $R$ is noetherian. In general, if $M\subseteq R^n$,
    then by (b), $M = M' + M_0$ where $M'$ is finitely generated
    and $M_0 = \ker(\pi_1|_{M})$ is isomorphic
    to a submodule of $R$, with embedding $M_0\to R^{n-1}$ given by
    $(0, r_2,\dots,r_n)\mapsto(r_2,\dots,r_n)$.
    By inductive hypothesis, $M_0$ is finitely generated.
    To finish, any $m\in M$ can be written as $m' + m_0$
    for some $m'\in M$ and $m_0\in M_0$, which are respectively
    $R$-linear combinations of the generators of $M'$ and $M_0$,
    so the generators of $M'$ and of $M_0$ together
    constitute a finite set of generators of $M$.
    \\

    (d) Let $M$ be a finitely generated $R$-module
    with generators $m_1$, \dots, $m_k$ and $N\subseteq M$ a submodule.
    Then, the morphism $\phi\colon R^k\to M$ given by
    $\phi(r_1,\dots,r_k) = r_1m_1 + \cdots + r_km_k$
    is surjective, so $M\cong R^k/\ker(\phi)$.
    Letting $\pi\colon R^k\to R^k/\ker(\phi)$ be the projection morphism,
    the submodule $\phi\inv(N)$ of $R^k/\ker(\phi)$ corresponding to $N$
    pulls back to a submodule $\pi\inv(\phi\inv(N))\subseteq R^k$,
    which by (c) is finitely generated, say with generators $n_1$, \dots, $n_\ell$.
    Then, if $n\in\pi\inv(\phi\inv(N))$ is an $R$-linear combination
    of the $n_j$, its image $n + \ker(\phi)$ under $\pi$
    descends to an $R$-linear combination of the $n_j + \ker(\phi)$,
    so $\phi\inv(N)$ is generated by the finite set
    $\{n_j + \ker(\phi) : j=1,\dots,\ell\}$, hence $N$ is finitely generated.
\end{soln}










\end{document}
